%% main.tex -- "Why Coding? A Guide for Families": parent information deck.
%%
%% Build from this folder with `make` (deck, notes and handout); see README.md.
%% Speaker-notes edition: notes.tex defines \kcwithnotes and inputs this file.
%% Text shared with the handout lives in content/shared.tex.
\documentclass[20pt,t]{beamer}
\usetheme{kidscode}
\usepackage[american]{babel}
\usepackage{csquotes}
\usepackage[style=apa,backend=biber,uniquename=false]{biblatex}
\addbibresource{references.bib}
\renewcommand*{\bibfont}{\footnotesize}
\usepackage{pgfplots}
\pgfplotsset{compat=1.18,
  % numbers in charts use the slide font, not the math font
  /pgf/number format/assume math mode=true}
\usepackage{array,booktabs,tabularx}
%% content/shared.tex -- text shared by the slide deck (main.tex) and the
%% one-page handout (handout.tex), so the two can never disagree.
%%
%% EDIT BEFORE PRESENTING: the three placeholders below.
\newcommand{\Presenter}{[Presenter name]}
\newcommand{\ContactLine}{[email or phone]}
\newcommand{\StartDate}{January 2027}  % sample calendar in each Teacher Guide

%% ---------------------------------------------------------------- course facts
%% Source of truth: grade-4-6/SYLLABUS.md, docs/DESIGN.md, docs/ASSUMPTIONS.md.
\newcommand{\CourseName}{Kids Coding}
\newcommand{\LevelWeeks}{10--12 weeks}   % per level; holiday weeks without a session do not count
\newcommand{\TotalWeeks}{30--36 weeks}   % all three levels
\newcommand{\CourseShape}{3 levels \textperiodcentered{} \LevelWeeks{} each}
\newcommand{\SessionLength}{one 60-minute session a week}
\newcommand{\GroupSize}{up to 4 children, working in 2 pairs}
\newcommand{\Homework}{15--20 minutes a week, optional but encouraged}
\newcommand{\Equipment}{a Mac, the free Scratch app and Thonny; no accounts, no internet during lessons}

%% Level summaries: \LevelXName, \LevelXTools, \LevelXLearn, \LevelXProject, \LevelXProjectLine
\newcommand{\LevelOneName}{Introductory}
\newcommand{\LevelOneTools}{Scratch blocks, then Python turtle}
\newcommand{\LevelOneLearn}{steps in order, loops, events and choices, variables, first Python}
\newcommand{\LevelOneProject}{My Interactive Art Show}
\newcommand{\LevelOneProjectLine}{a program that asks the viewer questions and draws a picture that changes with their answers}

\newcommand{\LevelTwoName}{Intermediate}
\newcommand{\LevelTwoTools}{Python}
\newcommand{\LevelTwoLearn}{functions, lists and dictionaries, files, testing, keyboard-driven games}
\newcommand{\LevelTwoProject}{Arcade Game or Art Studio}
\newcommand{\LevelTwoProjectLine}{an arcade game with score and high scores, or an art studio that saves designs, with the child's own tests}

\newcommand{\LevelThreeName}{Advanced}
\newcommand{\LevelThreeTools}{Python projects + git}
\newcommand{\LevelThreeLearn}{objects and classes, search and sort, recursion, data, how simple AI learns}
\newcommand{\LevelThreeProject}{Build Something Big}
\newcommand{\LevelThreeProjectLine}{a grid game, art app or data/AI app with classes, an algorithm, tests and a git history}

%% ---------------------------------------------------------------- how parents help
\newcommand{\HelpItems}{%
  \item \textbf{Ask for a demo:} \enquote{Show me what you made.}
  \item \textbf{Ask about the bug,} not the score: \enquote{What went wrong? How did you find it?}
  \item \textbf{Praise strategies and effort,} not \enquote{being a natural.}
  \item \textbf{Protect 15 minutes} a week for homework, at a regular time.
  \item \textbf{Encourage everyone equally:} girls hear less encouragement than boys.
  \item \textbf{Come to the showcase} at the end of each level.}

%% ---------------------------------------------------------------- safety promises
\newcommand{\SafetyItems}{%
  \item \textbf{No accounts, no logins, no email addresses.}
  \item \textbf{Nicknames only;} programs never ask for personal information.
  \item \textbf{Nothing is uploaded;} projects stay on the computers in the room.
  \item \textbf{Screen breaks every session,} plus unplugged activities.
  \item \textbf{Color-blind-safe materials} and a dyslexia-friendly edition.}


\ifdefined\kcwithnotes
  \setbeameroption{show notes}
\fi

\title{Why Coding? A Guide for Families}
\author{\Presenter}
\hypersetup{pdftitle={Why Coding? A Guide for Families},
  pdfauthor={Kids Coding},
  pdfsubject={Parent information evening for the Kids Coding course},
  pdflang={en-US}}

%% TikZ styles (defined here because '#' cannot appear inside a frame body)
\tikzset{
  box/.style={draw=#1, line width=1.4pt, fill=#1!10, rounded corners=6pt, align=center,
              minimum height=1.2in, text width=2.25in, font=\small},
  arr/.style={-{Stealth[length=10pt]}, line width=1.6pt, kcgray},
  card/.style={draw=#1, line width=1.3pt, rounded corners=8pt, fill=#1!8,
               minimum width=3.6in, minimum height=1.9in, anchor=north west},
  ladder/.style={draw=#1, line width=1.4pt, fill=#1!12, rounded corners=6pt,
               minimum width=1.9in, minimum height=0.7in, font=\semibold}}

%% Small helpers used on several slides
\newcommand{\bignum}[1]{{\fontsize{72}{72}\selectfont\extrabold\color{kcbluetext}#1}}
\newcommand{\figalt}[1]{\par\vskip0.04in{\scriptsize\color{kcgray}\textit{Figure:} #1\par}}
\newcommand{\gray}[1]{\textcolor{kcgray}{#1}}

\begin{document}

%% =====================================================================
%% 1. Title
%% =====================================================================
{\setbeamertemplate{footline}{}
\begin{frame}[plain]
  \begin{tikzpicture}[remember picture,overlay]
    \foreach \c [count=\i from 0] in {kcorange,kcsky,kcgreen,kcyellow,kcblue,kcvermillion,kcpurple}
      \fill[\c] ([xshift=\i*\paperwidth/7]current page.north west) rectangle ++(\paperwidth/7,-0.09in);
    % original illustration: a stack of code blocks that turns into lines of text code
    \begin{scope}[shift={([xshift=-5.3in,yshift=-1.9in]current page.north east)}, x=1in, y=1in]
      \foreach \c/\w/\y in {kcorange/1.35/0, kcblue/1.75/-0.72, kcgreen/1.5/-1.44, kcpurple/1.9/-2.16, kcsky/1.2/-2.88}
        \fill[\c, rounded corners=7pt] (0,\y) rectangle ++(\w,-0.52);
      \foreach \y in {0,-0.72,-1.44,-2.16,-2.88}
        \draw[-{Stealth[length=10pt]}, kcgray, line width=1.8pt] (2.1,\y-0.26) -- (2.5,\y-0.26);
      \foreach \w/\y/\x in {1.1/0/0, 1.4/-0.72/0.3, 0.9/-1.44/0.3, 1.3/-2.16/0.6, 0.8/-2.88/0}
        \fill[kcink!85, rounded corners=2pt] (2.75+\x,\y-0.19) rectangle ++(\w,-0.14);
      \node[anchor=north west, font=\footnotesize, text=kcgray] at (0,-3.55) {blocks};
      \node[anchor=north west, font=\footnotesize, text=kcgray] at (2.75,-3.55) {text code};
    \end{scope}
    \node[anchor=west, text width=7.0in, align=left]
      at ([xshift=0.8in,yshift=0.15in]current page.west) {%
      {\small\semibold\color{kcbluetext}\addfontfeatures{LetterSpace=8}FAMILY INFORMATION EVENING\par}
      \vskip0.2in
      {\fontsize{60}{66}\selectfont\bfseries Why Coding?\par}
      \vskip0.15in
      {\Large For families of Grade~5 learners\par}
      \vskip0.4in
      What the research says, how our course works, and how you can help at home.\par
      \vskip0.5in
      {\small\color{kcgray}\Presenter\enspace\textbar\enspace \CourseName{} course\enspace\textbar\enspace starts \StartDate\par}};
  \end{tikzpicture}
  \kcnotes{Welcome everyone and thank them for coming. Introduce yourself and your role in
    running the course. Set expectations: about 25 minutes of talk, then questions. Say up front
    that you will be honest about what research does and does not show: this is not a sales
    pitch, and nothing tonight promises jobs, grades or college places.}
  {The illustration is original: colored blocks (like Scratch) on the left turn into lines of
    typed code on the right. That is the journey the children take in this course.}
  {\qa{Do I need to know anything about computers to follow tonight?}{No. Every technical word
    is explained, and you never need to code to help your child.}}
\end{frame}}

%% =====================================================================
%% 2. Agenda
%% =====================================================================
\kicker{Tonight}
\begin{frame}{Tonight in 25 minutes}
  \vskip0.25in
  \begin{tikzpicture}[x=2.35in, y=1in]
    \foreach \n/\t/\s/\c [count=\i from 0] in
      {1/Why now/{Computing, AI and the skills ahead}/kcorange,
       2/The evidence/{What studies show, and what they don't}/kcsky,
       3/Good teaching/{What works, according to research}/kcgreen,
       4/Our course/{Three levels, one session a week}/kcblue,
       5/Your role/{Easy ways to help at home}/kcpurple}{
      \fill[\c] (\i,0) rectangle ++(0.92,0.1);
      \node[anchor=north west, inner sep=0pt, font=\fontsize{44}{44}\selectfont\extrabold, text=kcink!75] at (\i,-0.2) {\n};
      \node[anchor=north west, inner sep=0pt, text width=2.1in, font=\semibold] at (\i,-1.0) {\t};
      \node[anchor=north west, inner sep=0pt, text width=2.05in, font=\small, text=kcgray] at (\i,-1.45) {\s};
    }
  \end{tikzpicture}
  \vskip0.3in
  \begin{keybox}[Then]
    Questions, common concerns and next steps. Ask anything, any time.
  \end{keybox}
  \kcnotes{Walk through the five parts quickly. Invite people to interrupt with questions; a
    small group works best as a conversation. Mention that a one-page handout summarizes
    everything, so nobody needs to take notes.}
  {}
  {\qa{Will the slides be shared?}{Yes: hand out the one-page summary; the full slides with
    references are available on request.}}
\end{frame}

%% =====================================================================
%% 3. Why now: skills
%% =====================================================================
\kicker{Why now}
\begin{frame}{Computing is changing the skills our kids need}
  \vskip0.1in
  Employers expect these to be the \textbf{fastest-growing skills to 2030}:
  \vskip0.2in
  \begin{tikzpicture}[x=3.9in]
    \foreach \n/\t/\i/\c [count=\k from 0] in
      {1/AI and big data/\faRobot/kcblue,
       2/Networks and cybersecurity/\faShield*/kcgreen,
       3/Technological literacy/\faLaptopCode/kcvermillion}{
      \node[draw=\c, line width=1.5pt, fill=\c!10, rounded corners=8pt, minimum width=3.6in,
            minimum height=1.15in, anchor=north west] at (\k,0) {};
      \node[anchor=north west, font=\Large, text=\c!70!black] at (\k+0.04,-0.14) {\i};
      \node[anchor=north west, font=\semibold, text width=2.7in] at (\k+0.2,-0.17) {\n.\enspace\t};
    }
  \end{tikzpicture}
  \vskip0.25in
  \textbf{Analytical thinking} is the top \emph{core} skill: 7 in 10 employers call it essential.
  \begin{keybox}
    Not every child should become a programmer, but every child will live with this technology.
  \end{keybox}
  \source{\textcite{wef2025jobs}, a survey of over 1,000 employers.}
  \kcnotes{The World Economic Forum surveys large employers worldwide every two years. In the
    January 2025 report, the three skills expected to grow fastest by 2030 are AI and big data;
    networks and cybersecurity; and technological literacy. Analytical thinking remains the
    skill employers most often call essential (seven in ten). Stress the take-away: we are not
    trying to turn ten-year-olds into software engineers. We want them to understand, question
    and shape the technology they will grow up with.}
  {This is a survey of what employers \emph{expect}, not a measurement of the future. Forecasts
    like this are useful for direction, not for precise predictions.}
  {\qa{Will coding guarantee my child a good job?}{No, and nobody can promise that. What we can
    say is that understanding how computers and AI work is becoming a basic literacy, like
    understanding money or science.}}
\end{frame}

%% =====================================================================
%% 4. Why now: schools
%% =====================================================================
\begin{frame}{Schools are still catching up}
  \vskip0.2in
  \begin{columns}[T,onlytextwidth]
    \column{0.47\textwidth}
      \bignum{60\%}\par\vskip0.1in
      of US public high schools offered a foundational computer science class in 2024--25.
      Access has \textbf{stopped growing}.
    \column{0.47\textwidth}
      \bignum{6.1\%}\par\vskip0.1in
      of high school students were enrolled in one. Most students never take a class.
  \end{columns}
  \vskip0.3in
  \begin{keybox}
    A small, steady start outside school fills a gap that many schools cannot yet fill.
  \end{keybox}
  \source{\textcite{codeorg2025state}; national data, 2024--25 school year.}
  \kcnotes{These are the latest national figures from the annual State of Computer Science
    Education report (Code.org with the Computer Science Teachers Association and the ECEP
    Alliance), published December 2025. About 60 percent of public high schools offered a
    foundational computer science course, and about 6 percent of high school students were
    enrolled in one. The share of schools offering it has flattened compared with earlier years.
    The point is not that schools are failing; it is that most children will meet very little
    computer science in school, so a gentle start at home can help.}
  {These numbers describe US public high schools. Elementary schools are measured differently,
    so we do not quote a figure for them. Local schools may be better or worse than the
    national picture.}
  {\qa{Does our local school teach coding?}{It varies a lot by school. It is worth asking; our
    course works alongside school, not instead of it.}}
\end{frame}

%% =====================================================================
%% 5. Why now: AI
%% =====================================================================
\begin{frame}{With AI, understanding code matters more}
  \vskip0.1in
  \begin{columns}[T,onlytextwidth]
    \column{0.48\textwidth}
      \begin{cardbox}[kcblue]{\faFlask\enspace 69 learners aged 10--17}
        \begin{itemize}
          \item With an AI helper, they wrote more working code.
          \item They could still edit code by hand.
          \item A week later: no clear difference.
        \end{itemize}
      \end{cardbox}
    \column{0.48\textwidth}
      \begin{cardbox}[kcvermillion]{\faGlobeEurope\enspace OECD report, 2026}
        \begin{itemize}
          \item General chatbots can lift scores without lasting learning.
          \item Tools designed for learning work better.
        \end{itemize}
      \end{cardbox}
  \end{columns}
  \vskip0.12in
  \evOpen\enspace\gray{How AI will change learning to code is not yet known.}
  \begin{keybox}
    Someone still has to understand, check and fix what the computer produces.
  \end{keybox}
  \source{\textcite{kazemitabaar2023codex}; \textcite{oecd2026deo}.}
  \kcnotes{Many parents ask: if AI can write code, why learn? Two pieces of evidence. First, a
    controlled study at the University of Toronto with 69 learners aged 10 to 17: half had an AI
    code generator while learning Python. Those with AI completed more tasks and scored higher
    while learning, and they did not get worse at editing code by hand. A week later, on tests
    without AI, the difference was small and not statistically significant. Second, the OECD's
    2026 Digital Education Outlook warns that general-purpose chatbots can raise performance on
    tasks without producing lasting learning, while tools designed for learning (giving hints,
    asking questions) do better. Our course uses no online AI tools; in Level~3 the children
    build tiny AI models by hand so they understand how AI learns and where it fails.}
  {\enquote{Not statistically significant} means the difference was small enough that it could
    be down to chance. One study of 69 children cannot settle the question.}
  {\qa{Should my child use ChatGPT for homework?}{Not for this course. The homework is about
    thinking through the problem; we teach how AI works in Level 3 instead.}}
\end{frame}

%% =====================================================================
%% 6. What coding builds
%% =====================================================================
\kicker{What the research shows}
\begin{frame}{What coding builds, and how sure we are}
  \vskip0.1in
  \renewcommand{\arraystretch}{1.5}
  \begin{tabularx}{\textwidth}{@{}>{\semibold\RaggedRight}m{4.7in}>{\RaggedRight\arraybackslash}X>{\raggedleft\arraybackslash}m{2.75in}@{}}
    Coding and computational thinking & \gray{learns to program} & \evStrong\\
    Problem solving and planning & \gray{on coding-like tasks} & \evStrong\\
    Math, creativity, reasoning & \gray{smaller in strict studies} & \evMixed\\
    Reading and school grades & \gray{little effect} & \evMixed\\
    General persistence and teamwork & \gray{rarely measured} & \evOpen\\
  \end{tabularx}
  \begin{keybox}[What \enquote{computational thinking} means]
    Breaking a problem into steps, spotting patterns, and testing ideas until they work.
  \end{keybox}
  \source{\textcite{scherer2019transfer,scherer2020teaching}; \textcite{montuori2024ct}; \textcite{wing2006ct}.}
  \kcnotes{This slide is the honest summary of the whole evidence section. Read the rows from
    top to bottom: we are most confident about what is closest to coding itself. Children
    reliably learn to program and to think in steps. There is good evidence for gains in problem
    solving and planning, especially on tasks that resemble coding. Effects on math, creativity
    and reasoning are positive on average but shrink in the most careful studies, so we call
    them mixed. Reading and school grades benefit least. Whether coding builds general
    persistence or teamwork is an open question: few studies measure it well.}
  {Each badge uses a color, an icon and a word, so it reads the same in black and white.
    \enquote{Strong} means several careful studies and reviews agree. \enquote{Mixed} means
    results vary or shrink with stronger study designs. \enquote{Open} means we simply do not
    know yet.}
  {\qa{So will it help with math?}{Possibly a little, and especially with problem solving, but
    we do not promise better math grades. Think of it as its own valuable subject.}}
\end{frame}

%% =====================================================================
%% 7. Scherer meta-analysis
%% =====================================================================
\begin{frame}{The biggest review: 105 studies}
  \vskip0.1in
  \begin{itemize}
    \item Children reliably \textbf{learn to program} when taught.
    \item Some gains spread: \textbf{creative thinking, math, thinking about thinking}.
    \item Smaller gains for spatial skills and reasoning.
    \item Least benefit: \textbf{reading and school grades}.
    \item Stronger study designs found \textbf{smaller effects}.
  \end{itemize}
  \begin{keybox}
    Expect better coders and problem solvers, not automatically higher grades.
  \end{keybox}
  \source{\textcite{scherer2019transfer} (105 studies, 539 effects); \textcite{scherer2020teaching}.}
  \kcnotes{This is the largest review of whether learning to program helps other kinds of
    thinking. Scherer, Siddiq and S\'anchez Viveros combined 105 studies with 539 results. Skills
    closest to programming improved most. Gains also appeared for creative thinking, math and
    metacognition (thinking about your own thinking), and to a lesser degree spatial skills and
    reasoning. Reading and school achievement benefited least. Importantly, studies that
    compared coders with children who did nothing at all found bigger effects than studies with
    an active comparison group, and published studies showed bigger effects than unpublished
    ones. That is a reason to be modest. A companion review by the same team (2020, 139
    interventions) found a large effect of teaching on programming skill itself.}
  {A meta-analysis pools many studies to find the overall pattern. When effects get smaller in
    better-designed studies, part of the \enquote{benefit} in weaker studies may come from
    simply doing something new and fun, not from coding itself.}
  {\qa{What is \enquote{thinking about thinking}?}{Metacognition: noticing how you are solving a
    problem and changing strategy when it isn't working. Debugging practices exactly this.}}
\end{frame}

%% =====================================================================
%% 8. Chart: Montuori 2024
%% =====================================================================
\begin{frame}{Where coding seems to help most}
  \vskip0.05in
  \begin{columns}[T,onlytextwidth]
    \column{0.64\textwidth}
      \begin{tikzpicture}
        \begin{axis}[
          xbar, width=6.2in, height=3.4in,
          xmin=0, xmax=1.1, bar width=0.34in,
          symbolic y coords={Inhibition,Working memory,Planning,Problem solving},
          ytick=data, enlarge y limits=0.18,
          y tick label style={font=\small},
          xtick={0,0.2,0.4,0.6,0.8,1.0},
          x tick label style={font=\scriptsize, /pgf/number format/fixed}, xticklabel={\pgfmathprintnumber{\tick}},
          xlabel={Effect size (Hedges' \textit{g})}, xlabel style={font=\small},
          axis x line*=bottom, axis y line*=left,
          nodes near coords={\pgfmathprintnumber[fixed,precision=2,zerofill]{\pgfplotspointmeta}},
          nodes near coords style={font=\small\semibold, text=kcink},
          clip=false,
        ]
          \draw[kcgray, dashed, line width=0.8pt] ([yshift=-0.3in]axis cs:0.2,Inhibition) -- ([yshift=0.3in]axis cs:0.2,Problem solving) node[above, font=\scriptsize, text=kcgray] {small};
          \draw[kcgray, dashed, line width=0.8pt] ([yshift=-0.3in]axis cs:0.5,Inhibition) -- ([yshift=0.3in]axis cs:0.5,Problem solving) node[above, font=\scriptsize, text=kcgray] {medium};
          \draw[kcgray, dashed, line width=0.8pt] ([yshift=-0.3in]axis cs:0.8,Inhibition) -- ([yshift=0.3in]axis cs:0.8,Problem solving) node[above, font=\scriptsize, text=kcgray] {large};
          \addplot[fill=kcblue, draw=none] coordinates
            {(0.17,Inhibition) (0.20,Working memory) (0.36,Planning) (0.89,Problem solving)};
        \end{axis}
      \end{tikzpicture}
      \figalt{bar chart of effect sizes: problem solving 0.89, planning 0.36, working memory
        0.20, inhibition 0.17. Dashed lines mark the usual small, medium and large benchmarks.}
    \column{0.32\textwidth}
      \vskip0.15in
      \textbf{19 experiments} with children; 11 pooled.
      \vskip0.25in
      Big gain in \textbf{problem solving}; small gains in self-control and memory.
      \vskip0.25in
      \evMixed\par\vskip0.05in
      \gray{Few studies; effects vary.}
  \end{columns}
  \source{\textcite{montuori2024ct}; pooled (fixed-effect) estimates from 11 studies.}
  \kcnotes{This 2024 review looked only at experiments with children: 19 studies with 1,523
    participants, of which 11 could be combined. The pooled effects were large for problem
    solving (0.89), smaller for planning (0.36), and small for working memory (0.20) and
    inhibition, the ability to stop an automatic response (0.17). The dashed lines show common
    rules of thumb: 0.2 is small, 0.5 medium, 0.8 large.}
  {An effect size compares the average improvement of the coding group with the comparison
    group, in units of how spread out children's scores are. 0.89 means a big difference; 0.17
    a small one. Because the review pooled only 11 studies, we label it mixed rather than
    strong.}
  {\qa{Is 0.89 a lot?}{Yes, it is large by usual standards, but it comes from a small number of
    studies, so treat it as promising rather than proven.}}
\end{frame}

%% =====================================================================
%% 9. Arfé RCT
%% =====================================================================
\begin{frame}{A closer look: one fair test}
  \vskip0.3in
  \begin{tikzpicture}
    \node[box=kcgray] (kids) {\textbf{179 first graders}\\5 schools,\\no coding before};
    \node[box=kcgray, right=0.42in of kids] (rand) {\textbf{Randomly split}\\by class,\\like a coin toss};
    \node[box=kcblue, right=0.42in of rand, yshift=0.75in] (code) {\textbf{Coding}\\8 hours\\over 4 weeks};
    \node[box=kcorange, right=0.42in of rand, yshift=-0.75in] (ctrl) {\textbf{Usual STEM}\\\textbf{lessons}\\same time};
    \node[box=kcgreen, right=0.42in of code, yshift=-0.75in] (res) {\textbf{Result}\\coding group improved more in \textbf{planning}};
    \draw[arr] (kids) -- (rand);
    \draw[arr] (rand.east) -- ++(0.2,0) |- (code.west);
    \draw[arr] (rand.east) -- ++(0.2,0) |- (ctrl.west);
    \draw[arr] (code.east) -- ++(0.2,0) |- (res.west);
    \draw[arr] (ctrl.east) -- ++(0.2,0) |- (res.west);
  \end{tikzpicture}
  \figalt{flow diagram of a randomized trial: 179 first graders randomly assigned by class to
    coding or usual STEM lessons; the coding group improved more in planning.}
  \vskip0.25in
  \textbf{Limits:} young children, a short program, one country (Italy). Gains in self-control were smaller.
  \source{\textcite{arfe2020coding}, a cluster-randomized controlled trial.}
  \kcnotes{This is one of the cleanest studies in the field. Researchers in Italy randomly
    assigned ten first-grade classes either to eight hours of coding over four weeks (using
    Code.org puzzles) or to their usual STEM lessons. Before and after, children did planning
    puzzles such as the Tower of London and tests of inhibition, the ability to stop an automatic
    response. The coding group improved more, especially in planning; the effect on inhibition
    was smaller.}
  {Random assignment is the gold standard: because chance decides who codes, differences at the
    end are much more likely to be caused by the coding itself. Still, it is one short study of
    six-year-olds, so we should not assume the same size of effect for our ten-year-olds.}
  {\qa{Why show a study of six-year-olds?}{Because it is one of the few truly randomized
    trials. Studies of older children mostly point the same way but with weaker designs.}}
\end{frame}

%% =====================================================================
%% 10. Early experience
%% =====================================================================
\begin{frame}{Early experiences shape interest}
  \vskip0.1in
  \begin{columns}[T,onlytextwidth]
    \column{0.48\textwidth}
      \begin{cardbox}[kcgreen]{\faRobot\enspace Programming a robot, age 6}
        Afterwards, girls were as interested and confident as boys.\par
        \vskip0.05in\gray{Stereotypes did not change.}
      \end{cardbox}
    \column{0.48\textwidth}
      \begin{cardbox}[kcorange]{\faChartLine\enspace Plans at age 13--14}
        Expecting a science career predicted a science degree (2--3$\times$).\par
        \vskip0.05in\gray{A link, not proof of cause.}
      \end{cardbox}
  \end{columns}
  \vskip0.1in
  Girls report less encouragement in computing from parents and teachers.
  \begin{keybox}
    Early, positive experiences keep doors open. They do not decide a future.
  \end{keybox}
  \source{\textcite{master2017programming}; \textcite{tai2006planning}; \textcite{google2016diversity}.}
  \kcnotes{Three findings about interest. First, in a University of Washington experiment,
    first graders who programmed a small robot showed no gender gap in technology interest and
    confidence afterwards, while in the comparison groups boys reported more of both. The short
    activity did not change children's stereotypes, though. Second, a large US study followed
    students from eighth grade (1988) into college: those who expected a science-related career
    at age 13 to 14 were roughly two to three times as likely to graduate with a science degree.
    This is a correlation (an association), not proof that early interest causes the degree.
    Third, a 2016 Google--Gallup survey found girls are less likely than boys to be told by
    parents and teachers that they would be good at computer science.}
  {Correlation versus causation: children who already love science might both expect a science
    career and later study it. The early expectation might not cause the degree; it might just
    be an early sign of it. The 2006 study is also old (data from 1988 to 2000).}
  {\qa{My daughter says coding is for boys. What can I do?}{Encourage her as much as you would a
    son, and point out the variety of people who code. Our course never assigns roles by gender.}}
\end{frame}

%% =====================================================================
%% 11. What we don't know
%% =====================================================================
\begin{frame}{What we don't know (yet)}
  \vskip0.1in
  \begin{itemize}
    \item Gains in other school subjects are \textbf{mixed and often small}.
    \item Stronger studies find \textbf{smaller benefits}.
    \item Skills rarely carry far: the same is true of \textbf{chess and music}.
    \item Kids who choose coding may differ: \textbf{correlation is not cause}.
    \item Results depend heavily on \textbf{how it is taught}.
    \item We make \textbf{no promises} about jobs, pay or college.
  \end{itemize}
  \begin{keybox}[Our promise]
    Capable, confident beginner programmers who practice problem solving every week.
  \end{keybox}
  \source{\textcite{pea1984cognitive} (an early critique); \textcite{sala2017far}; \textcite{scherer2019transfer}.}
  \kcnotes{This is the most important slide for trust. Researchers have argued about this since
    the 1980s: Pea and Kurland at Bank Street College questioned the hope that programming
    automatically improves general thinking, and early studies with the Logo language were
    disappointing when there was no deliberate teaching. (That critique is old, which we flag,
    but its lesson still holds.) A 2017 review of chess, music and memory training found the same
    general pattern: skills rarely carry over to unrelated areas, and apparent effects shrink in
    better-designed studies. So our promise is modest: children will become capable, confident
    beginner programmers and practice problem solving. Anything beyond that is a bonus.}
  {\enquote{Transfer} means using a skill learned in one area somewhere else. Near transfer (to
    similar problems) is common; far transfer (to unrelated subjects) is rare for any activity.}
  {\qa{Then why bother?}{Because coding is worth learning for its own sake, like music or a
    language, and because the way we teach it (projects, teamwork, debugging) builds good
    learning habits.}}
\end{frame}

%% =====================================================================
%% 12. Six principles
%% =====================================================================
\kicker{What good teaching looks like}
\begin{frame}{Six principles the evidence supports}
  \vskip0.15in
  \begin{tikzpicture}[x=3.9in, y=-2.1in]
    \foreach \t/\d/\s/\i/\c [count=\k from 0] in
      {Real projects/{Games and art they love}/{Resnick et al., 2009}/\faPaintBrush/kcorange,
       {Low floor, high ceiling}/{Easy start, room to grow}/{Resnick et al., 2009}/\faLayerGroup/kcsky,
       Code in pairs/{Take turns at the keys}/{Umapathy \& Ritzhaupt, 2017}/\faUserFriends/kcgreen,
       Teach debugging/{A step-by-step bug hunt}/{Michaeli \& Romeike, 2019}/\faBug/kcvermillion,
       {Off-screen, read first}/{Predict, then run}/{Sentance et al., 2019}/\faEye/kcblue,
       {Spiral, with feedback}/{Ideas return, deeper}/{Bruner, 1960}/\faRedo/kcpurple}{
      \pgfmathtruncatemacro{\col}{mod(\k,3)}
      \pgfmathtruncatemacro{\row}{\k/3}
      \node[card=\c] (c\k) at (\col,\row) {};
      \node[anchor=north west, font=\Large, text=\c!70!black] at ([xshift=0.18in,yshift=-0.16in]c\k.north west) {\i};
      \node[anchor=north west, font=\semibold, text width=3.25in] at ([xshift=0.18in,yshift=-0.66in]c\k.north west) {\t};
      \node[anchor=north west, text width=3.25in] at ([xshift=0.18in,yshift=-1.05in]c\k.north west) {\d};
      \node[anchor=south east, font=\scriptsize, text=kcgray] at ([xshift=-0.15in,yshift=0.1in]c\k.south east) {\s};
    }
  \end{tikzpicture}
  \kcnotes{Research is clearer about \emph{how} to teach coding than about its spill-over
    benefits. These six principles come up again and again. The next three slides give the
    evidence for each, and then we show where each one appears in our course. The small label in
    the corner of each card names the key source.}
  {}
  {\qa{Is there one best way to teach coding?}{No single method wins everywhere, but these
    ingredients are consistently linked to better learning and motivation.}}
\end{frame}

%% =====================================================================
%% 13. Create, don't consume
%% =====================================================================
\begin{frame}{Create, don't just consume}
  \begin{columns}[T,onlytextwidth]
    \column{0.54\textwidth}
      \begin{itemize}
        \item Kids learn best making things \textbf{they care about}.
        \item \textbf{Low floor:} a first program in minutes.
        \item \textbf{High ceiling:} real complexity later.
        \item \textbf{Wide walls:} games, art, stories, data.
      \end{itemize}
    \column{0.42\textwidth}
      \begin{tikzpicture}
        \node[ladder=kcsky] (u) at (0,0) {Use};
        \node[ladder=kcblue] (m) at (1.15in,0.8in) {Modify};
        \node[ladder=kcbluetext] (c) at (2.3in,1.6in) {Create};
        \draw[arr] (u.east) -| (m.south);
        \draw[arr] (m.east) -| (c.south);
        \node[font=\scriptsize, text=kcgray, anchor=north] at (u.south) {play with a program};
        \node[font=\scriptsize, text=kcgray, anchor=south] at (c.north) {build your own};
      \end{tikzpicture}
      \figalt{a staircase from Use to Modify to Create.}
  \end{columns}
  \vskip-0.1in
  \begin{keybox}
    Every level ends with a project the child designs, builds and presents.
  \end{keybox}
  \source{\textcite{papert1980mindstorms}; \textcite{resnick2009scratch}; \textcite{lee2011ct}.}
  \kcnotes{Seymour Papert at MIT argued in 1980 that children learn best by building things that
    matter to them. The Scratch team (Resnick and colleagues, 2009) turned this into three design
    goals: a low floor (easy to start), a high ceiling (room for complex projects) and wide walls
    (many kinds of project, so every child finds an interest). Lee and colleagues (2011) describe
    a progression called Use--Modify--Create: first play with a working program, then change it,
    then design your own. Our lessons follow exactly that pattern.}
  {These are design principles from experienced researchers rather than results of one
    experiment. They are widely used in computing education because they keep children engaged.}
  {\qa{My child only wants to make games. Is that OK?}{Yes. Games are a great vehicle: they need
    loops, conditions, variables and debugging. Art and data projects are there for others.}}
\end{frame}

%% =====================================================================
%% 14. Pairs, debugging, feedback
%% =====================================================================
\begin{frame}{Learn together, fix bugs on purpose}
  \vskip0.1in
  \begin{columns}[T,onlytextwidth]
    \column{0.32\textwidth}
      \begin{cardbox}[kcgreen]{\faUserFriends\enspace Pairs}
        Better assignments, exams and persistence. Beginners gain most.
      \end{cardbox}
    \column{0.32\textwidth}
      \begin{cardbox}[kcvermillion]{\faBug\enspace Debugging}
        A taught, step-by-step process improved skill \textbf{and} confidence.
      \end{cardbox}
    \column{0.32\textwidth}
      \begin{cardbox}[kcpurple]{\faComments\enspace Feedback}
        Helps most when it says \textbf{how to improve}.
      \end{cardbox}
  \end{columns}
  \vskip0.2in
  \evStrong\enspace Pairs and feedback.\par\vskip0.1in
  \evMixed\enspace Debugging: promising, but few studies.
  \source{\textcite{umapathy2017pair}; \textcite{denner2014pair}; \textcite{michaeli2019debugging}; \textcite{wisniewski2020feedback}.}
  \kcnotes{Pair programming means two children share one computer: the driver types, the
    navigator reads, plans and spots mistakes, and they swap regularly. A 2017 review of 18
    studies (mostly university courses) found better results for pairs on assignments, exams and
    staying in the course, though not on attitudes. A study of 320 middle-school students
    (Denner and colleagues, 2014) found pairs gained more programming knowledge, especially
    children with less experience. Debugging: a German classroom study found that explicitly
    teaching a systematic process for finding bugs improved both debugging performance and
    students' confidence compared with practice alone. Feedback: a 2020 review of 435 studies
    found a solid medium-sized effect on learning, largest when feedback gives information about
    the task and how to improve, and smaller for praise alone.}
  {You may hear a famous figure of 0.79 for feedback (Hattie and Timperley, 2007). A later,
    larger re-analysis by the same research group found about 0.48, which is still a solid
    effect. We quote the more careful estimate qualitatively.}
  {\qa{What if my child does not get along with their partner?}{We rotate pairs, set clear roles
    and swap every few minutes. With an odd number, a trio uses driver, navigator and tester.}}
\end{frame}

%% =====================================================================
%% 15. Unplugged, read first, spiral
%% =====================================================================
\begin{frame}{Off-screen first, read before writing, return often}
  \vskip0.1in
  \begin{columns}[T,onlytextwidth]
    \column{0.32\textwidth}
      \begin{cardbox}[kcblue]{\faCut\enspace Unplugged}
        Starting off-screen helped young pupils' thinking and motivation.
      \end{cardbox}
    \column{0.32\textwidth}
      \begin{cardbox}[kcsky]{\faEye\enspace Predict, then run}
        Students aged 11--14 who read code first beat a comparison group.
      \end{cardbox}
    \column{0.32\textwidth}
      \begin{cardbox}[kcpurple]{\faRedo\enspace Spiral}
        Big ideas come back each level, deeper each time.
      \end{cardbox}
  \end{columns}
  \vskip0.25in
  \textbf{Loops, for example:} Scratch \emph{repeat} $\rightarrow$ Python \texttt{for}
    $\rightarrow$ game loops $\rightarrow$ recursion.
  \source{\textcite{delolmo2020unplugged}; \textcite{sentance2019primm}; \textcite{bruner1960process}.}
  \kcnotes{Unplugged activities teach computing ideas without a computer: acting out a loop as a
    dance, sorting cards, giving a parent-robot exact instructions. A Spanish study with early
    primary pupils found that starting with unplugged activities before moving to the computer
    helped computational thinking and motivation. PRIMM (Predict, Run, Investigate, Modify, Make)
    asks students to read and predict what code does before writing their own; an evaluation in
    13 UK schools with 493 students aged 11 to 14 found better post-test results than a
    comparison group. The spiral curriculum idea comes from psychologist Jerome Bruner (1960):
    return to big ideas repeatedly, at greater depth each time.}
  {Reading code before writing it is like reading books before writing stories: it builds a
    picture of how programs work, so writing becomes less frustrating.}
  {\qa{Isn't time away from the computer wasted?}{No. Off-screen activities make the ideas
    concrete and give eyes a break; they are short and always lead into building.}}
\end{frame}

%% =====================================================================
%% 16. How our course applies the principles
%% =====================================================================
\kicker{Our course}
\begin{frame}{How our course puts this into practice}
  \vskip0.05in
  \renewcommand{\arraystretch}{1.18}
  \begin{tabularx}{\textwidth}{@{}>{\semibold\RaggedRight}p{3.6in}>{\RaggedRight\arraybackslash}X@{}}
    \toprule
    Principle & In our course\\
    \midrule
    Real projects & A final project every level; 10+ free-build times\\
    Low floor, high ceiling & Scratch, then Python; easier and harder paths\\
    Code in pairs & Driver and navigator, swapped every session\\
    Teach debugging & A \enquote{bug hunt} every week, with a strategy\\
    Off-screen, read first & 6+ unplugged activities; predict-then-run\\
    Spiral, with feedback & Ideas return each level; exit check each lesson\\
    \bottomrule
  \end{tabularx}
  \source{Course design (docs/DESIGN.md, docs/SYLLABUS.md); per-level minimums are checked automatically.}
  \kcnotes{This slide links research to practice. Every one of the 42 sessions follows the same
    five-part shape. The minimums in this table (per level: at least six unplugged activities,
    ten bug hunts, ten pair-programming activities and ten free-build times) are checked
    automatically by the course's quality tests, so they cannot quietly disappear when material is
    edited. Each lesson includes a support path for children who need more time and an extension
    path for those who finish early. Many lessons ask children to predict what code will do
    before running it.}
  {An \enquote{exit check} is a quick question at the end of a session that shows whether the
    idea landed, so the next session can adjust.}
  {\qa{What if my child is ahead or behind the others?}{Every lesson has an easier route and a
    stretch challenge, and two spare sessions per level give time to catch up.}}
\end{frame}

%% =====================================================================
%% 17. Blocks then Python
%% =====================================================================
\begin{frame}{Why blocks first, then Python}
  \vskip0.1in
  \begin{columns}[T,onlytextwidth]
    \column{0.5\textwidth}
      \begin{itemize}
        \item Blocks remove typing errors: kids focus on ideas.
        \item Teens learned more with blocks at first.
        \item After switching to text, the gap disappeared.
        \item So: 6 weeks of Scratch, then Python.
      \end{itemize}
    \column{0.46\textwidth}
      \begin{tikzpicture}
        \begin{axis}[
          width=5.4in, height=2.4in, xmin=-0.3, xmax=0.8, ymin=0.4, ymax=1.6,
          ytick=\empty, axis y line=none, axis x line*=bottom,
          xtick={-0.2,0,0.2,0.4,0.6}, x tick label style={font=\scriptsize, /pgf/number format/fixed}, xticklabel={\pgfmathprintnumber{\tick}},
          xlabel={Effect on learning (Hedges' \textit{g})}, xlabel style={font=\small},
          clip=false]
          \draw[kcgray, dashed, line width=1pt] (axis cs:0,0.4) -- (axis cs:0,1.6)
            node[above, font=\scriptsize, text=kcgray] {no difference};
          \draw[kcblue, line width=4pt] (axis cs:-0.078,1) -- (axis cs:0.567,1);
          \fill[kcbluetext] (axis cs:0.245,1) circle (7pt);
          \node[below=6pt, font=\small\semibold] at (axis cs:0.245,1) {0.25};
          \node[anchor=east, font=\scriptsize, text=kcgray] at (axis cs:0.8,1.45) {favors blocks $\rightarrow$};
        \end{axis}
      \end{tikzpicture}
      \figalt{estimate 0.25 with a range from \textminus0.08 to 0.57 that crosses zero: a small
        edge for blocks that could be chance.}
  \end{columns}
  \vskip0.1in
  \evMixed\enspace Review of 13 studies: a small edge for blocks, could be chance.
  \source{\textcite{weintrop2017blocks,weintrop2019transition}; \textcite{xu2019blocks}.}
  \kcnotes{Scratch, created at MIT, uses snap-together blocks, so children cannot make spelling
    or bracket mistakes. In a high school study by Weintrop and Wilensky, students using a
    block-based tool scored higher on concept tests after the first five weeks than students
    using an equivalent text tool. When both groups then moved to a professional text language
    (Java), the difference disappeared: blocks gave no harm and no lasting advantage. A 2019
    review of 13 studies found a small advantage for blocks that was not statistically
    significant. Our choice: blocks for the first six weeks while ideas like loops and events are
    new, then Python, a real language used by professionals, with every Scratch idea re-taught in
    text.}
  {The chart shows one estimate (the dot) and its range of plausible values (the line). Because
    the line crosses zero, the true effect could be zero. That is why we call it mixed, and why
    we switch to text quickly instead of staying with blocks for a long time.}
  {\qa{Isn't Scratch just a toy?}{It is a real programming language with the same core ideas; we
    use it as a ramp, and move to Python after six weeks.}}
\end{frame}

%% =====================================================================
%% 18. Roadmap timeline
%% =====================================================================
\begin{frame}{The roadmap: three levels, about a year}
  \vskip0.3in
  \begin{tikzpicture}[x=0.27in, y=1in]
    \foreach \lv/\name/\tools/\months/\c/\x in
      {1/Introduction/{Scratch, then Python turtle}/{Jan--Apr}/kcbluetext/0,
       2/Intermediate/{Python: functions, data, games}/{May--Aug}/kcgreentext/14.6,
       3/Advanced/{Python projects + git}/{Sep--Dec}/kcvermilliontext/29.2}{
      \node[anchor=south west, font=\small, text=kcgray] at (\x,0.72) {\months};
      \fill[\c, rounded corners=5pt] (\x,0) rectangle ++(14,0.7);
      \node[anchor=west, text=white, font=\semibold] at (\x+0.35,0.35) {Level \lv: \name};
      \node[anchor=north west, font=\small, text width=3.6in] at (\x,-0.1) {\tools};
      \foreach \w/\ic in {6/\faClipboardCheck, 10.5/\faTools, 12.8/\faStar}{
        \fill[\c] (\x+\w,-0.72) circle (0.17in);
        \node[text=white, font=\scriptsize] at (\x+\w,-0.72) {\ic};
      }
      \draw[\c, line width=1.2pt] (\x+0.2,-0.72) -- (\x+13.8,-0.72);
    }
    \node[anchor=north west, font=\small, text=kcgray] at (0,-1.1)
      {\faClipboardCheck\ checkpoint, week 6\qquad \faTools\ final project, weeks 10--11\qquad
       \faStar\ showcase, week 12\qquad (sample calendar)};
  \end{tikzpicture}
  \figalt{timeline of three equal levels across one year: Level 1 Introduction (January to
    April), Level 2 Intermediate (May to August), Level 3 Advanced (September to December), each
    with a checkpoint, a final project and a showcase.}
  \begin{keybox}
    No experience needed: everyone starts at Level 1, week 1.
  \end{keybox}
  \kcnotes{Walk along the timeline. Each level is twelve weekly sessions plus two spare sessions
    for holidays, absences or extra practice, so about four months. With the sample calendar,
    Level 1 runs January to April, Level 2 May to August with a short summer gap, and Level 3
    September to December. Inside every level the rhythm is the same: new ideas in weeks 1 to 5,
    a review and friendly checkpoint in week 6, more new ideas in weeks 7 to 9, a final project in
    weeks 10 and 11, and a showcase in week 12. Families can decide level by level whether to
    continue.}
  {}
  {\qa{Do we have to commit to all three levels?}{No. Each level ends with a complete project and
    a certificate; you decide about the next level then.}}
\end{frame}

%% =====================================================================
%% 19-21. Levels
%% =====================================================================
%% \levelslide{no.}{name}{color}{tools line}{learn items}{final project}{project line}{examples}{progress}
\newcommand{\levelslide}[9]{%
  \begin{frame}{Level #1: #2}
    \vskip-0.08in
    {\small\semibold\color{#3}#4\par}
    \vskip0.15in
    \begin{columns}[T,onlytextwidth]
      \column{0.5\textwidth}
        {\semibold Your child will learn to\dots\par}
        \vskip0.05in
        \begin{itemize}#5\end{itemize}
      \column{0.46\textwidth}
        \begin{cardbox}[#3]{\faStar\enspace Final project}
          \textbf{#6}\par\vskip0.03in #7
        \end{cardbox}
        \vskip0.08in
        \textbf{Along the way:} #8\par
    \end{columns}
    \vskip0.12in
    \textbf{Progress checks:} #9\par
  }

\levelslide{1}{\LevelOneName}{kcbluetext}{\LevelOneTools{} \textperiodcentered{} weeks 1--6 in blocks, 7--12 in Python}
  {\item give exact step-by-step instructions
   \item repeat with loops; react with \enquote{if}
   \item keep score with variables
   \item find and fix bugs with a strategy
   \item write first Python programs}
  {\LevelOneProject}{Asks the viewer questions, then draws a scene from the answers.}
  {catch games, star patterns, a magic dice.}
  {welcome survey (ungraded), week 6, week 12, project rubric.}
  \source{Level 1 syllabus; aligned with CSTA standards for grades 3--5 \parencite{csta2017standards}.}
  \kcnotes{Level 1 starts from zero. Weeks 1 to 6 use Scratch blocks: making sprites move, draw
    patterns, react to keys, and keep score in simple games. Week 6 is review and a friendly
    checkpoint. Weeks 7 to 12 move to Python, using the turtle, which draws on screen: the same
    ideas, now typed. The final project, My Interactive Art Show, must use a question, a variable,
    an if, a loop, something random and a function the child wrote. Children who need more time
    can build it in Scratch with the same rubric.}
  {A \enquote{variable} is a labeled box that remembers something, like a score. A
    \enquote{function} is a named mini-program you can reuse, like a recipe card.}
  {\qa{Is typing a problem for a ten-year-old?}{Programs at this level are short; we start very
    small, and typing improves quickly.}}
\end{frame}

\levelslide{2}{\LevelTwoName}{kcgreentext}{\LevelTwoTools{} in Thonny, a free beginner editor}
  {\item organize code into reusable functions
   \item store data in lists and dictionaries
   \item save and load files, like high scores
   \item write tests that check code automatically
   \item build keyboard-controlled games}
  {\LevelTwoProject}{A game with saved high scores, or an art studio, plus the child's own tests.}
  {secret-code machines, pixel art, guess-the-number.}
  {placement quiz, week 6, week 12, project rubric.}
  \source{Level 2 syllabus; aligned with CSTA standards for grades 6--8 \parencite{csta2017standards}.}
  \kcnotes{Level 2 is all Python. Children learn to organize code into functions that return
    answers, store many values in lists and dictionaries, save data to files, and, a highlight,
    write automatic tests that prove their code works. They build keyboard-driven games and
    generative art. A short placement quiz in week 1 helps tailor support; it is never a pass or
    fail. The final project is a choice between an arcade game and an art studio, with at least
    three tests the child writes.}
  {A \enquote{test} is a tiny program that checks another program gives the right answer, like a
    spell-checker for logic. Professionals use them every day.}
  {\qa{What if my child struggled with Level 1?}{The placement quiz and Level 1 review show where
    to start; repeating parts of Level 1 is completely fine.}}
\end{frame}

\levelslide{3}{\LevelThreeName}{kcvermilliontext}{\LevelThreeTools{} \textperiodcentered{} git is a \enquote{save-game} system for code}
  {\item design programs with classes and objects
   \item compare ways to search and sort
   \item draw fractals with recursion
   \item clean and chart a small data set
   \item explain how simple AI learns, and fails}
  {\LevelThreeProject}{A game, art app or AI app with classes, an algorithm and tests.}
  {virtual pets, fractal trees, a tiny AI.}
  {placement quiz, week 6, week 12, project rubric.}
  \source{Level 3 syllabus; AI topics follow the AI4K12 Five Big Ideas \parencite{touretzky2019ai4k12}.}
  \kcnotes{Level 3 introduces the ideas behind larger programs: classes and objects (blueprints
    and things built from them), searching and sorting, recursion for fractal art, analyzing a
    small made-up data set, and how simple AI learns from examples, and copies their mistakes.
    The children build a tiny AI (a nearest-neighbor classifier) by hand rather than using online
    AI services, so no data leaves the computer. They also use git, the tool professionals use to
    save versions of their code; the final project needs at least five tests and eight meaningful
    git commits. The AI topics follow the AI4K12 framework from AAAI and CSTA; some Level 3 goals
    borrow from the grades 9--10 standards on purpose.}
  {A \enquote{nearest-neighbor classifier} guesses what something is by finding the most similar
    example it has seen, which shows children exactly why biased examples give biased answers.}
  {\qa{Isn't this too advanced for 11-year-olds?}{It is challenging by design, with step-by-step
    support. By Level 3 the children have about eight months of practice behind them.}}
\end{frame}

%% =====================================================================
%% 22. A typical session
%% =====================================================================
\begin{frame}{What a 60-minute session looks like}
  \vskip0.15in
  \begin{tikzpicture}[x=0.195in]
    \foreach \m/\name/\c/\x/\ic/\fg in
      {8/Warm-up/kcorange/0/\faRunning/kcink,
       10/New idea/kcsky/8/\faLightbulb/kcink,
       10/Demo/kcblue/18/\faChalkboardTeacher/white,
       24/Build in pairs/kcgreen/28/\faUserFriends/kcink,
       8/Show \& tell/kcpurple/52/\faStar/kcink}{
      \fill[\c, rounded corners=4pt] (\x+0.08,0) rectangle (\x+\m-0.08,1.0in);
      \node[font=\large, text=\fg] at (\x+\m/2,0.66in) {\ic};
      \node[font=\small\semibold, text=\fg] at (\x+\m/2,0.27in) {\m\,min};
      \node[anchor=north, font=\small, align=center, text width=\m*0.19in] at (\x+\m/2,-0.08in) {\name};
    }
  \end{tikzpicture}
  \figalt{five-part session: warm-up 8 minutes, new idea 10, demo 10, building in pairs 24, show
    and tell 8.}
  \vskip0.15in
  \begin{columns}[T,onlytextwidth]
    \column{0.31\textwidth}
      \begin{cardbox}[kcblue]{\faHome\enspace Homework}
        15--20 min a week, optional
      \end{cardbox}
    \column{0.31\textwidth}
      \begin{cardbox}[kcgreen]{\faLaptop\enspace Equipment}
        A Mac with free apps; no accounts
      \end{cardbox}
    \column{0.31\textwidth}
      \begin{cardbox}[kcpurple]{\faUsers\enspace Group}
        Up to 4 kids in 2 pairs
      \end{cardbox}
  \end{columns}
  \kcnotes{Every session has the same five parts, so children always know what is coming: a
    warm-up game (often unplugged), one new idea with a kid-friendly analogy, a short live demo,
    about 24 minutes building in pairs with a few minutes of free building, and a wrap-up where
    children show their work and answer a quick exit question. Homework takes 15 to 20 minutes a
    week, with a core task for everyone and optional challenge and creative extras; none needs
    the internet. Equipment: a Mac with the free Scratch app and Thonny. In Level 2 a testing tool
    called pytest is installed once (internet needed once, no account), and Level 3 uses Apple's
    free command-line tools for git.}
  {}
  {\qa{We don't have a Mac at home. Is that a problem?}{No. Paper parts of the homework work at
    the kitchen table, and computer tasks can be finished at the next session.}}
\end{frame}

%% =====================================================================
%% 23. What success looks like
%% =====================================================================
\begin{frame}{What success looks like}
  \vskip0.05in
  \renewcommand{\arraystretch}{1.25}
  \begin{tabularx}{\textwidth}{@{}>{\semibold}p{2.2in}>{\RaggedRight\arraybackslash}X@{}}
    \toprule
    By the end of & most children can\dots\\
    \midrule
    \textcolor{kcbluetext}{Level 1} & build a small game and a drawing program, and explain them\\
    \textcolor{kcgreentext}{Level 2} & organize a game into functions, save scores, test code\\
    \textcolor{kcvermilliontext}{Level 3} & plan a multi-part project, track it in git, explain choices\\
    \bottomrule
  \end{tabularx}
  \vskip0.15in
  \textbf{How we check:} checkpoint in week 6, exit check in week 12, project rubric.\par
  \vskip0.05in
  \textbf{Rubric:} Just starting \textperiodcentered{} Getting there \textperiodcentered{} Got it \textperiodcentered{} Expert
  \begin{keybox}[No grades]
    Scores guide support. Creativity, working code and explaining it count equally.
  \end{keybox}
  \kcnotes{These are realistic outcomes, not guarantees: children progress at different rates.
    Each level has a short, friendly checkpoint in week 6 and an exit check in week 12, plus a
    project rubric with four levels. The rubric values three things equally: making the project
    their own, making it work, and being able to explain it. Scores are used to decide what
    support to offer next, never as grades.}
  {A \enquote{rubric} is a simple chart that describes what each level of work looks like, so
    children can see what \enquote{better} means and assess themselves.}
  {\qa{Will I get a report?}{You will see the project at each showcase and can ask for the
    checkpoint and rubric results; there is also a progress tracker children fill in.}}
\end{frame}

%% =====================================================================
%% 24. How parents can help
%% =====================================================================
\kicker{Your role}
\begin{frame}{How you can help, no coding needed}
  \vskip0.1in
  \begin{itemize}
    \HelpItems
  \end{itemize}
  \source{\textcite{google2016diversity}; \textcite{master2017programming}; \textcite{wisniewski2020feedback}.}
  \kcnotes{None of this needs any coding knowledge. Asking for a demo makes your child the
    expert, which builds confidence. Asking about bugs signals that mistakes are normal and
    interesting. Praising strategies (\enquote{you tried three ways to fix that}) gives more useful
    feedback than labels like \enquote{you're a natural}. A regular 15-minute slot is easier than a
    long weekend session. Encouraging daughters and sons equally matters because surveys show
    girls hear less encouragement about computing.}
  {Feedback research shows that specific comments about the task and strategy help learning more
    than general praise.}
  {\qa{What if my child asks me for help and I don't know the answer?}{Perfect: say \enquote{Let's
    find out.} Ask them to explain the code line by line; they often spot the bug themselves.}}
\end{frame}

%% =====================================================================
%% 25. Safety and inclusion
%% =====================================================================
\kicker{Safety and inclusion}
\begin{frame}{Safe, private and inclusive by design}
  \vskip0.1in
  \begin{columns}[T,onlytextwidth]
    \column{0.56\textwidth}
      \begin{itemize}
        \SafetyItems
      \end{itemize}
    \column{0.4\textwidth}
      \begin{cardbox}[kcgreen]{\faShield*\enspace Why no accounts?}
        US law requires parental consent before online services collect data from children
        under 13. We collect nothing.
      \end{cardbox}
  \end{columns}
  \source{\textcite{ftc2025coppa}; \textcite{aap2016media}; \textcite{okabe2008cud}.}
  \kcnotes{Privacy: the course uses only offline apps, so no child creates an account, enters an
    email or uploads anything. Programs ask for nicknames or favorite colors, never real names or
    addresses. This is deliberate: the federal COPPA rule (updated in 2025) requires online
    services to get verifiable parental consent before collecting personal information from
    children under 13, and we prefer not to collect anything at all. Screen time: the American
    Academy of Pediatrics recommends consistent family limits for school-age children that
    protect sleep, physical activity and family time; our sessions include a stretch break, the
    20-20-20 eye rule and unplugged activities. Inclusion: all materials use the Okabe--Ito
    color-blind-safe palette with icons and words alongside color, a highly legible font designed
    by the Braille Institute, and a dyslexia-friendly edition on request.}
  {The 2016 pediatric guidance does not set a single number of hours for this age; it asks
    families to make a plan. One hour of creative coding a week plus 15 minutes of homework is a
    small part of a week.}
  {\qa{Will my child's photo or work be posted online?}{No. Showcases happen on screen at home,
    and nothing is uploaded.}}
\end{frame}

%% =====================================================================
%% 26. Common questions
%% =====================================================================
\kicker{Questions}
\begin{frame}{Common questions}
  \vskip0.1in
  \begin{columns}[T,onlytextwidth]
    \column{0.48\textwidth}
      \begin{cardbox}[kcsky]{Isn't 10 too young?}
        No. Scratch was designed for this age.
      \end{cardbox}
      \vskip0.1in
      \begin{cardbox}[kcsky]{Won't AI make this pointless?}
        Someone has to understand and check AI's work.
      \end{cardbox}
    \column{0.48\textwidth}
      \begin{cardbox}[kcsky]{More screen time?}
        One hour a week, with breaks and games.
      \end{cardbox}
      \vskip0.1in
      \begin{cardbox}[kcsky]{I can't code. Can I help?}
        Yes: ask questions and praise effort.
      \end{cardbox}
  \end{columns}
  \kcnotes{These are the four questions parents ask most often. Keep answers short and invite
    follow-ups. If a question comes up that the evidence cannot answer, say so: \enquote{Research
    doesn't tell us that yet.}}
  {}
  {\qa{What if my child loses interest?}{It happens. Children choose their projects, which helps,
    and each level is a natural stopping point.}
   \qa{Is this better than an online course?}{Different: a small group, a familiar adult, no
    accounts, and a lot of talking and pair work, which research supports.}}
\end{frame}

%% =====================================================================
%% 27. Next steps
%% =====================================================================
\begin{frame}{Next steps}
  \vskip0.15in
  \begin{enumerate}
    \item \textbf{Take the handout} home and talk it over with your child.
    \item \textbf{Let us know} if your child will join Level 1 (starts \StartDate).
    \item \textbf{Check the computer:} we help install Scratch and Thonny.
  \end{enumerate}
  \vskip0.3in
  \begin{keybox}[Questions?]
    \Presenter\enspace\textbar\enspace\ContactLine
  \end{keybox}
  \kcnotes{Close by thanking parents again. Summarize in one sentence: coding is worth learning
    for its own sake; the research gives us good reasons to teach it through projects, pairs and
    debugging; and we are honest about what it will and will not do. Then open the floor.}
  {}
  {\qa{Can siblings join?}{The group is up to four children so everyone gets attention; ask, and
    we will see what is possible.}}
\end{frame}

%% =====================================================================
%% Appendix
%% =====================================================================
\appendix
\sectionslide{A}{Appendix}{Evidence summary, notes on sources, credits and references}

\kicker{Appendix}
\begin{frame}{Evidence at a glance}
  \renewcommand{\arraystretch}{1.12}
  \footnotesize
  \begin{tabularx}{\textwidth}{@{}>{\RaggedRight}X>{\RaggedRight}p{2.6in}>{\RaggedRight\arraybackslash}p{4.0in}@{}}
    \toprule
    Claim & Strength & Key sources\\
    \midrule
    Children learn to program when taught & \evStrong & \textcite{scherer2020teaching}\\
    Gains in problem solving and planning & \evStrong & \textcite{montuori2024ct}; \textcite{arfe2020coding}\\
    Gains in math, creativity, reasoning & \evMixed & \textcite{scherer2019transfer}\\
    Gains in reading and school grades & \evMixed & \textcite{scherer2019transfer}\\
    Pair programming improves learning & \evStrong & \textcite{umapathy2017pair}; \textcite{denner2014pair}\\
    Blocks before text help learning & \evMixed & \textcite{xu2019blocks}; \textcite{weintrop2019transition}\\
    Explicit debugging teaching helps & \evMixed & \textcite{michaeli2019debugging}\\
    Coding builds general persistence & \evOpen & no strong studies found\\
    Long-term effect of AI on learning to code & \evOpen & \textcite{kazemitabaar2023codex}; \textcite{oecd2026deo}\\
    \bottomrule
  \end{tabularx}
  \kcnotes{Use this slide if parents ask for the evidence behind a specific claim. The problem
    solving row is marked strong for coding-like tasks; the Montuori chart on slide 8 is marked
    mixed because that single review pooled only 11 studies.}{}{}
\end{frame}

\begin{frame}{Notes on dates and limits}
  \small
  \begin{itemize}
    \item \textbf{Older sources:} \textcite{pea1984cognitive}, \textcite{papert1980mindstorms} and
      \textcite{bruner1960process} are cited for ideas, not current data.
    \item \textbf{\textcite{tai2006planning}} uses data from 1988--2000; it shows an association only.
    \item \textbf{Most pair-programming studies} were with university students.
    \item \textbf{Hattie \& Timperley (2007)} reported 0.79 for feedback; a larger re-analysis
      found about 0.48 \parencite{wisniewski2020feedback}.
    \item \textbf{CSTA} published a draft revision of its standards in 2025; the 2017 version is cited.
    \item \textbf{Survey data} (WEF, Code.org, Google--Gallup) describe expectations and access, not outcomes.
  \end{itemize}
  \kcnotes{This slide documents where sources are old or limited, so the presenter can answer
    critical questions honestly.}{}{}
\end{frame}

\begin{frame}{Credits and licenses}
  \small
  \begin{itemize}
    \item \textbf{All diagrams and charts:} original, drawn in TikZ/pgfplots for this deck.
    \item \textbf{Chart data:} published values only, from the cited papers.
    \item \textbf{Fonts:} Atkinson Hyperlegible Next and Mono (Braille Institute), SIL Open Font License 1.1.
    \item \textbf{Icons:} Font Awesome 5 Free via the fontawesome5 package; icons CC BY 4.0, font SIL OFL 1.1.
    \item \textbf{Colors:} Okabe--Ito color-blind-safe palette \parencite{okabe2008cud}.
    \item \textbf{Trademarks:} Scratch is a project of the Scratch Foundation; Python is a trademark of the PSF.
  \end{itemize}
  \kcnotes{No photographs or third-party images are used, so there are no image licenses to track.}{}{}
\end{frame}

\begin{frame}[allowframebreaks]{References}
  \printbibliography[heading=none]
  \kcnotes{Full reference list generated from references.bib.}{}{}
\end{frame}

\end{document}
